\begin{lemma}
\label{editorial-source-lemma-presentation-sym-exterior}
Let $R$ be a ring.
Let $M_2 \to M_1 \to M \to 0$ be an exact sequence of $R$-modules.
For every integer $n \geq 1$ there are exact sequences
$$
M_2 \otimes_R \text{Sym}^{n - 1}(M_1)
\to
\text{Sym}^n(M_1)
\to
\text{Sym}^n(M)
\to
0
$$
and similarly
$$
M_2 \otimes_R \wedge^{n - 1}(M_1)
\to
\wedge^n(M_1)
\to
\wedge^n(M)
\to
0
$$
In degree zero both induced maps are the identity map $R \to R$.
\end{lemma}

\begin{proof}
Denote the maps by $u : M_2 \to M_1$ and $v : M_1 \to M$.
We prove both assertions, writing $C(M_1)$ for the symmetric algebra
or the exterior algebra in the respective case. Let $J$ be the
homogeneous two sided ideal of $C(M_1)$ generated by all
$u(y)$ in degree one, for $y \in M_2$.
The map $v$ induces a graded $R$-algebra map
$C(M_1)/J \to C(M)$. Conversely, the assignment
$v(x) \mapsto x \bmod J$ defines an $R$-linear map
$M \to (C(M_1)/J)^1$: it is independent of the lift $x$
because $\Ker(v) = \Im(u)$. Sending a pure tensor in $M$
to the product of these images defines a map from its tensor algebra,
with the identity map on the degree-zero part $R$.
In the symmetric case it kills the commutator relations.
In the exterior case it kills $m\otimes m$, since a lift $x$ of $m$
satisfies $x\wedge x = 0$ in $C(M_1)$.
Thus it induces a graded $R$-algebra map
$C(M) \to C(M_1)/J$. The two maps are inverse on $R$ and on
the degree-one generators, and hence on the entire algebras.

\medskip\noindent
It follows that the kernel of $C^n(M_1) \to C^n(M)$ is $J^n$.
For $n \geq 1$ every homogeneous term generating $J^n$ has the form
$a\,u(y)\,b$, with $a$ of degree $p$, $b$ of degree $q$ and
$p+1+q=n$. In the symmetric algebra this equals $u(y)ab$.
In the exterior algebra it equals $(-1)^p u(y)ab$:
on a pure degree-$p$ tensor this is obtained by $p$ successive
interchanges of degree-one factors, and the assertion extends by
$R$-linearity. Therefore $J^n$ is exactly the image of the map
$$
M_2\otimes_R C^{n-1}(M_1) \longrightarrow C^n(M_1),
\qquad y\otimes z \longmapsto u(y)z.
$$
The reverse inclusion holds because every $u(y)$ belongs to $J$.
This proves the two displayed exact sequences, including their
surjective last maps. In degree zero $J^0=0$, and both maps
$C^0(M_1)=R \to C^0(M)=R$ are the identity.
\end{proof}